\subsection{Continuity and failure of differentiability}

A function may be continuous at a point and still fail to have a derivative there.

\begin{theorembox}[Differentiability implies continuity]
If \(f\) is differentiable at \(a\), then \(f\) is continuous at \(a\).
\end{theorembox}

The converse is false.

\begin{examplebox}[A corner]
For
\[
f(x)=|x|,
\]
the function is continuous at \(0\), but the left-hand slope is \(-1\) and the right-hand slope is \(1\). Since the one-sided difference quotients do not agree, \(f'(0)\) does not exist.
\end{examplebox}

\begin{examplebox}[A vertical tangent]
For
\[
f(x)=x^{1/3},
\]
the difference quotient near \(0\) becomes unbounded. The graph has a vertical tangent rather than a finite tangent slope, so the ordinary derivative at \(0\) is not a real number.
\end{examplebox}

\begin{interpretationbox}[What can prevent a derivative]
Non-differentiability may come from a discontinuity, a corner, a cusp, or an infinite slope. These are geometrically different failures, and the graph should be read before a formula is applied mechanically.
\end{interpretationbox}
